\documentclass[11pt,a4paper]{article}
\usepackage[margin=1in]{geometry}
\usepackage{times}
\usepackage{latexsym}
\usepackage{graphicx}
\usepackage{amsmath}
\usepackage{amssymb}
\usepackage{booktabs}
\usepackage{multirow}
\usepackage{xcolor}
\usepackage{subcaption}
\usepackage{natbib}
\usepackage[colorlinks=true,citecolor=blue,linkcolor=blue]{hyperref}

\title{Differential Depth Response in Non-Paraphrastic LLM Generation Chains: \\
An Empirical Characterization Across Model Families}

\author{Anonymous}

\begin{document}
\maketitle

%==============================================================================
% ABSTRACT
%==============================================================================
\begin{abstract}
Web-scraped corpora increasingly contain LLM-generated text of unknown provenance depth---the number of recursive generation steps from the original human-written source.
Current detection methods offer only binary real-vs-synthetic classification, conflating qualitatively different material.
We formulate \emph{generation depth classification} (gen-0/gen-1/gen-2+) for non-paraphrastic continuation chains and characterize a phenomenon we term \emph{differential depth response} (DDR): central-tendency features (entropy mean, TTR, hapax ratio) drift monotonically with depth ($|\rho| > 0.60$), while shape features (entropy max, tail decay) remain near-stationary ($|\rho| < 0.10$).
Across 8 base-model generators from 5 families (0.6B--8B), a lightweight 7-feature ensemble achieves leak-free binary F1 = 0.694 (95\% CI [0.688, 0.700]) under leave-one-generator-out (LOGO) evaluation, and 0.830 on cross-model chains using a single 410M reference model.
Per-generator binary F1 ranges from 0.633 to 0.724.
Depth-aware data curation is statistically equivalent to random subsampling for downstream causal language modeling quality (TOST $p = 1.1 \times 10^{-4}$ at $\pm$2\% margin), establishing depth fingerprinting as a diagnostic tool for provenance characterization rather than a curative tool for data selection.
The signal degrades substantially under instruction tuning and high temperature ($t \geq 0.9$), restricting current scope to base models.
We also identify \texttt{entropy\_std} as an intermediate feature whose classification is inconsistent across reference models, motivating a spectrum model (drifting $\to$ intermediate $\to$ stationary) over the original binary partition.
\end{abstract}

%==============================================================================
% 1. INTRODUCTION
%==============================================================================
\section{Introduction}

Web-scraped training data increasingly contains LLM-generated text of unknown provenance: recent estimates suggest 7--17\% of scientific papers may contain AI-generated content.
This recursive contamination creates \emph{generation chains} where output from one model becomes input to another, potentially degrading data quality across training iterations.
Current detection methods treat text classification as a binary problem---human-written vs.\ machine-generated---conflating gen-1 text (which retains much of the source distribution's lexical properties) with deeper generations that exhibit measurably reduced vocabulary diversity and compressed perplexity.

This binary framing is insufficient for data governance.
A dataset curator deciding whether to include questionable text needs to know not just \emph{whether} it was machine-generated, but \emph{how far removed} it is from the original human source.
We formulate \textbf{generation depth classification} as an ordinal prediction task: given a text passage, determine whether it is gen-0 (human-written), gen-1 (first-generation synthetic), or gen-2+ (recursively generated).

Our central empirical finding is \textbf{differential depth response (DDR)}: distributional features computed from a reference model's token-level negative log-likelihood (NLL) distribution exhibit systematically different behavior across generation depths.
Central-tendency and diversity features---entropy mean, type-token ratio (TTR), and hapax ratio---drift monotonically with depth, consistent with recursive generation inducing a location shift in the NLL distribution.
Shape features---entropy max and tail decay---remain approximately stationary, preserving distributional shape even as the center shifts.
This differential response is not an artifact of any particular reference model: it holds across three independent reference models from different families.

We validate DDR across 8 base-model generators from 5 families (Pythia, OLMo, Qwen, Llama, Mistral; 0.6B--8B) and present three core contributions:

\begin{enumerate}
    \item \textbf{Task formulation and empirical validation}: Generation depth classification beyond binary detection, with a 7-feature ensemble achieving leak-free binary F1 = 0.694 under 8-generator LOGO evaluation.
    \item \textbf{DDR characterization}: Drifting features show $|\rho| > 0.60$ while shape features show $|\rho| < 0.10$, with a $\sim$10$\times$ gap in KS statistics. We identify \texttt{entropy\_std} as an intermediate feature (Cohen's $d = 0.458$, TOST rejection rate 0/1000) that challenges a strict binary partition.
    \item \textbf{Negative finding on downstream utility}: Depth-aware data curation is equivalent to random subsampling for downstream CLM quality (TOST confirmed at $\pm$2\% margin), establishing depth fingerprinting as diagnostic, not curative.
\end{enumerate}

The signal degrades substantially under instruction tuning (OLMo-Instruct F1 = 0.338 vs.\ 0.763 base) and high sampling temperature ($t \geq 0.9$), restricting current scope to base models at moderate temperatures.
We report these boundaries transparently as defining constraints rather than limitations to be worked around.

%==============================================================================
% 2. RELATED WORK
%==============================================================================
\section{Related Work}

\paragraph{Perplexity-based detection.}
Statistical detection methods leverage the observation that machine-generated text tends to occupy high-probability regions of the source model's distribution.
DetectGPT \cite{mitchell2023detectgpt} uses perturbation-based curvature estimation; Fast-DetectGPT \cite{bao2024fast} replaces perturbations with conditional probability curvature.
GLTR \cite{gehrmann2019gltr} visualizes token-level log-probabilities.
Binoculars \cite{hans2024binoculars} compares cross-entropy ratios between two models.
DNA-GPT \cite{yang2024dnagpt} uses divergent n-gram analysis.
These methods provide binary classification and do not address ordinal depth estimation.
When repurposed as depth predictors, their continuous scores show moderate signal but high cross-generator variance.

\paragraph{Model collapse and iterative generation.}
Shumailov et al.\ study model collapse under recursive training on synthetic data, showing progressive loss of distributional tails.
The ``LLM as Broken Telephone'' framework (ACL 2025) examines iterative generation chains and bounds paraphrasing chains to $\sim$3 distinguishable classes under semantic-preservation constraints.
Our work differs by addressing non-paraphrastic continuation chains where semantic drift is unconstrained.

\paragraph{Data quality and vocabulary diversity.}
Vocabulary diversity metrics (TTR, hapax legomena ratio) have long been used in stylometry and authorship attribution.
Recent work connects these metrics to LLM-generated text detection, observing that machine-generated text exhibits reduced lexical diversity.
Our contribution situates these observations within an ordinal depth framework, showing that diversity reduction is progressive and measurable across generation steps.

\paragraph{OOD detection and membership inference.}
Our feature set overlaps with out-of-distribution detection approaches that use reference model statistics.
We explicitly acknowledge this connection: the contribution is not the features themselves but (a) the ordinal depth estimation formulation, (b) the empirical DDR characterization showing systematic drifting-vs-stationary behavior, and (c) the downstream utility analysis.

%==============================================================================
% 3. METHOD
%==============================================================================
\section{Method}

\subsection{Problem Formulation}
Given a text passage $x$, we aim to classify its \emph{generation depth} $d \in \{0, 1, 2+\}$, where $d=0$ denotes human-written text, $d=1$ denotes first-generation LLM output, and $d=2+$ denotes text produced by two or more recursive generation steps.
We construct \emph{generation chains} by iteratively feeding each generation's output as the prompt for the next, using non-paraphrastic continuation prompts.

\subsection{Feature Design}
We compute 7 features from the token-level NLL distribution obtained by a reference model $M_\text{ref}$ (Pythia-410M by default).
For each text $x = (x_1, \ldots, x_T)$, we compute $\ell_t = -\log p_{M_\text{ref}}(x_t | x_{<t})$ and derive:

\paragraph{Drifting features} (central-tendency/diversity):
\begin{itemize}
    \item \textbf{entropy\_mean}: $\bar{\ell} = \frac{1}{T}\sum_t \ell_t$
    \item \textbf{TTR}: type-token ratio $|\text{types}(x)| / T$
    \item \textbf{hapax\_ratio}: fraction of tokens appearing exactly once
\end{itemize}

\paragraph{Shape features} (near-stationary):
\begin{itemize}
    \item \textbf{entropy\_max}: $\max_t \ell_t$
    \item \textbf{tail\_decay}: exponential decay rate of the upper tail of $\{\ell_t\}$
\end{itemize}

\paragraph{Intermediate features}:
\begin{itemize}
    \item \textbf{entropy\_min}: $\min_t \ell_t$ (Spearman $\rho = -0.30$)
    \item \textbf{entropy\_std}: $\text{std}(\{\ell_t\})$ (contested; $\rho = -0.20$, Cohen's $d = 0.458$)
\end{itemize}

We classify features into a spectrum---drifting ($|\rho| > 0.50$), intermediate ($0.15 \leq |\rho| \leq 0.50$), and stationary ($|\rho| < 0.15$)---rather than a binary partition, based on TOST bootstrap validation (\S\ref{sec:tost}).

\subsection{Classifier and Evaluation Protocol}
We train an XGBoost classifier on the 7-feature representation.
To ensure leak-free evaluation, we use two protocols:

\paragraph{LOGO (Leave-One-Generator-Out):} Train on chains from all generators except one; test on the held-out generator. Repeat for each generator and average.

\paragraph{LOFO (Leave-One-Family-Out):} Same as LOGO but holding out an entire model family (e.g., all Pythia models), testing cross-family transfer.

Our \textbf{primary metric} is \emph{leak-free binary F1}: the F1 score for distinguishing gen-1 from gen-2+, excluding gen-0 samples to avoid seed-sharing leakage between generators that share the same human-written seeds.

%==============================================================================
% 4. EXPERIMENTS
%==============================================================================
\section{Experiments}

\subsection{Feature Ablation and DDR Trend Analysis}
\label{sec:ablation}

We first validate the feature design on MVP data using 5-fold cross-validation.
Table~\ref{tab:ablation} shows that the full 7-feature set (macro F1 = 0.820) significantly outperforms both surface-only (TTR + hapax; F1 = 0.763) and reference-only (entropy + tail; F1 = 0.719) subsets (McNemar $p \approx 0$ for both comparisons).

\begin{table}[t]
\centering
\small
\begin{tabular}{lccc}
\toprule
\textbf{Feature Set} & \textbf{Macro F1} & \textbf{95\% CI} & \textbf{$N$} \\
\midrule
Full (7-feature) & 0.820 & [0.805, 0.835] & 7 \\
Surface-only & 0.763 & [0.747, 0.774] & 2 \\
Reference-only & 0.719 & [0.706, 0.733] & 5 \\
\bottomrule
\end{tabular}
\caption{Feature ablation under 5-fold CV on MVP data. Surface features (TTR, hapax) and reference features (entropy, tail) provide complementary signal (McNemar $p \approx 0$).}
\label{tab:ablation}
\end{table}

Table~\ref{tab:trend} shows the DDR trend classification.
Drifting features exhibit strong monotonic relationships with depth ($|\rho| > 0.60$, KS $> 0.50$), while shape features are near-stationary ($|\rho| < 0.10$, KS $< 0.06$).
Notably, \texttt{entropy\_std} ($\rho = -0.20$, KS $= 0.15$, Cohen's $d = 0.458$) falls between groups, motivating the intermediate category.

\begin{table}[t]
\centering
\small
\begin{tabular}{lccl}
\toprule
\textbf{Feature} & \textbf{Spearman $\rho$} & \textbf{KS} & \textbf{Class} \\
\midrule
hapax\_ratio & $-0.63$ & 0.537 & Drifting \\
TTR & $-0.61$ & 0.521 & Drifting \\
entropy\_mean & $-0.60$ & 0.516 & Drifting \\
entropy\_min & $-0.30$ & 0.280 & Intermediate \\
entropy\_std & $-0.20$ & 0.150 & Intermediate \\
entropy\_max & $+0.05$ & 0.052 & Stationary \\
tail\_decay & $+0.04$ & 0.043 & Stationary \\
\bottomrule
\end{tabular}
\caption{Feature trend classification across depth. Drifting features show $\sim$10$\times$ larger KS statistics than stationary features. \texttt{entropy\_std} is classified as intermediate (TOST $d = 0.458$, rejection 0/1000).}
\label{tab:trend}
\end{table}

Perplexity validation with Pythia-410M confirms monotonic decrease across depth: median perplexity drops from 13.50 (gen-0) through 4.96, 3.00, 2.07, to 1.72 (gen-4).

\subsection{Cross-Generator Evaluation}
\label{sec:cross_gen}

We expand to 5 generators (Pythia-2.8B, Pythia-1B, OLMo-2-1B, Qwen3-1.7B, Qwen3-0.6B; $N = 23{,}032$ records) and evaluate under LOGO and LOFO protocols.

\begin{table}[t]
\centering
\small
\begin{tabular}{llccc}
\toprule
\textbf{Generator} & \textbf{Params} & \textbf{Bin. F1} & \textbf{Mac. F1} & \textbf{CI (Bin.)} \\
\midrule
Pythia-2.8B & 2.8B & 0.724 & 0.851 & [0.710, 0.738] \\
Pythia-1B & 1B & 0.718 & 0.831 & [0.703, 0.733] \\
Qwen3-1.7B & 1.7B & 0.702 & 0.805 & [0.687, 0.717] \\
OLMo-2-1B & 1.5B & 0.668 & 0.752 & [0.652, 0.684] \\
Qwen3-0.6B & 0.6B & 0.659 & 0.780 & [0.643, 0.675] \\
\midrule
\textbf{Average} & — & \textbf{0.694} & \textbf{0.804} & [0.688, 0.700] \\
\bottomrule
\end{tabular}
\caption{LOGO per-generator results (5-gen). Binary F1 ranges 0.659--0.724 across generators.}
\label{tab:logo}
\end{table}

LOFO evaluation yields binary F1 = 0.686 and macro F1 = 0.791, only 0.8 pp below LOGO (Table~\ref{tab:lofo}), indicating that within-family information leakage is minimal.

\begin{table}[t]
\centering
\small
\begin{tabular}{lcc}
\toprule
\textbf{Protocol} & \textbf{Binary F1} & \textbf{Macro F1} \\
\midrule
LOGO & 0.694 & 0.804 \\
LOFO & 0.686 & 0.791 \\
$\Delta$ & $-0.008$ & $-0.013$ \\
\bottomrule
\end{tabular}
\caption{LOGO vs.\ LOFO comparison. The small gap confirms limited family leakage.}
\label{tab:lofo}
\end{table}

\subsection{Cross-Reference-Model Validation}

To confirm that DDR is generation-intrinsic rather than a Pythia-specific artifact, we repeat the analysis with OLMo-2-1B and Qwen3-0.6B as alternative reference models.

\begin{table}[t]
\centering
\small
\begin{tabular}{lcccc}
\toprule
\textbf{Ref.\ Model} & \textbf{Params} & \textbf{LOGO F1} & \textbf{$\rho_\text{ent.mean}$} & \textbf{$\rho_\text{tail}$} \\
\midrule
Pythia-410M & 410M & 0.803 & $-0.60$ & $+0.04$ \\
OLMo-2-1B & 1.5B & 0.807 & $-0.52$ & $+0.06$ \\
Qwen3-0.6B & 0.6B & 0.816 & $-0.58$ & $+0.03$ \\
\bottomrule
\end{tabular}
\caption{Cross-reference-model validation. DDR trends hold across all reference models, confirming generation-intrinsic origin. Drift magnitude varies (OLMo range [$-0.61$, $-0.91$] vs.\ Pythia [$-1.19$, $-1.40$]).}
\label{tab:cross_ref}
\end{table}

All three reference models produce consistent DDR: entropy drift is universally negative, and tail\_decay remains near-stationary (Table~\ref{tab:cross_ref}).
However, \texttt{entropy\_std} shows inconsistent classification across references (Cohen's $d$: 0.458 Pythia, 0.310 OLMo, 0.520 Qwen), further supporting its intermediate status.

\subsection{Baseline Comparison}
\label{sec:baseline}

We compare the 7-feature ensemble against existing detection methods repurposed as ordinal depth predictors.

\begin{table}[t]
\centering
\small
\begin{tabular}{lccc}
\toprule
\textbf{Method} & \textbf{LOGO F1} & \textbf{95\% CI} & \textbf{Cost} \\
\midrule
7-feature XGBoost & 0.798 & [0.786, 0.799] & 1$\times$ \\
Avg-LogProb & 0.765 & [0.616, 0.820] & 1$\times$ \\
LD-Score & 0.721 & [0.588, 0.770] & 0$\times$ \\
FastDetectGPT & 0.627 & [0.511, 0.762] & 2$\times$ \\
Binoculars & 0.579 & [0.405, 0.664] & 2$\times$ \\
4-det + surface & 0.795 & [0.780, 0.810] & 4$\times$+ \\
\bottomrule
\end{tabular}
\caption{Baseline comparison (LOGO macro F1). The 7-feature ensemble significantly outperforms all individual detectors (bootstrap $p < 0.001$). The multi-detector ensemble matches performance at $>$2.4$\times$ inference cost.}
\label{tab:baseline}
\end{table}

The 7-feature ensemble significantly outperforms all individual baselines (bootstrap $p < 0.001$; Table~\ref{tab:baseline}).
Avg-LogProb is the strongest single detector (0.765) but with much wider confidence intervals, indicating higher cross-generator variance.
A multi-detector ensemble (4-det + surface features; 0.795) matches the 7-feature set but requires $>$2.4$\times$ inference cost and 3.2$\times$ GPU memory.
The 7-feature ensemble's advantage lies in efficiency and interpretability rather than absolute performance superiority.

\subsection{Cross-Model Chains}

Real-world web contamination involves multiple LLM sources.
We construct cross-model chains where each depth step uses a different generator from a 5-model pool, evaluating under generator-level LOGO.

\begin{table}[t]
\centering
\small
\begin{tabular}{lccc}
\toprule
\textbf{Setup} & \textbf{Bin. F1} & \textbf{3-Class F1} & \textbf{Protocol} \\
\midrule
Single-model (LOGO) & 0.694 & 0.762 & LOGO \\
Cross-model (LOGO) & 0.830 & 0.789 & Gen-LOGO \\
$\Delta$ & $+$0.136 & $+$0.027 & — \\
\bottomrule
\end{tabular}
\caption{Single-model vs.\ cross-model chains. Generator heterogeneity amplifies depth signal by $+$13.6 pp binary F1.}
\label{tab:cross_model}
\end{table}

Cross-model chains yield substantially higher binary F1 (0.830 vs.\ 0.694; Table~\ref{tab:cross_model}).
Generator heterogeneity amplifies depth-to-depth feature variance, making classification easier.
This is practically relevant: real-world multi-source contamination is the condition under which DDR is strongest.
Per-depth F1 shows gen-0 at 0.95, gen-1 at 0.82, and gen-2+ at 0.78.

\subsection{Downstream Utility (Negative Finding)}

We fine-tune OLMo-1B on depth-stratified subsets (3 seeds, bootstrap CIs) to test whether depth-aware curation improves downstream causal language modeling.

\begin{table}[t]
\centering
\small
\begin{tabular}{lccc}
\toprule
\textbf{Variant} & \textbf{$N$} & \textbf{PPL} & \textbf{$\Delta$ vs.\ SM} \\
\midrule
All depths & 12,299 & 2.335 & $-$0.032 \\
Size-matched random & 3,951 & 2.367 & — \\
Depth 0+1 only & 3,951 & 2.387 & $+$0.020 \\
Binary-filtered & 3,951 & 2.389 & $+$0.022 \\
Depth 0 only & 797 & 2.681 & $+$0.314 \\
\bottomrule
\end{tabular}
\caption{Downstream CLM perplexity by training data composition. At equal sample count, depth-aware curation does not improve over random subsampling.}
\label{tab:downstream}
\end{table}

At equal sample count ($n = 3{,}951$), size-matched random subsampling (PPL = 2.367) outperforms depth 0+1 filtering (2.387) and binary filtering (2.389; Table~\ref{tab:downstream}).
TOST equivalence is confirmed at $\pm$2\% margin ($p = 1.1 \times 10^{-4}$).
Seed variance is extremely small ($< 0.002$ PPL across seeds), indicating the result is robust.
This establishes depth fingerprinting as a \textbf{diagnostic} tool for provenance characterization, not a curative tool for data selection.

\subsection{9-Generator Scale Evaluation}
\label{sec:9gen}

Expanding to 9 generators (adding Gemma-2-2B, Llama-3.2-3B, Qwen3-8B, and Mistral-7B) tests broader generalization.

\begin{table}[t]
\centering
\small
\begin{tabular}{llccc}
\toprule
\textbf{Generator} & \textbf{Params} & \textbf{3-Cl F1} & \textbf{Bin F1} & \textbf{G1 F1} \\
\midrule
Pythia-1B & 1B & 0.857 & 0.718 & 0.65 \\
Qwen3-0.6B & 0.6B & 0.819 & 0.659 & 0.61 \\
Pythia-2.8B & 2.8B & 0.812 & 0.724 & 0.68 \\
Qwen3-1.7B & 1.7B & 0.811 & 0.702 & 0.66 \\
OLMo-2-1B & 1.5B & 0.773 & 0.668 & 0.58 \\
Llama-3.2-3B & 3B & 0.759 & 0.690 & 0.60 \\
Qwen3-8B & 8B & 0.696 & 0.647 & 0.52 \\
\textbf{Gemma-2-2B} & 2B & \textbf{0.628} & \textbf{0.440} & \textbf{0.036} \\
Mistral-7B & 7B & 0.605 & 0.633 & 0.44 \\
\bottomrule
\end{tabular}
\caption{9-generator LOGO results. Gemma-2-2B exhibits anomalous gen-1 F1 = 0.036; Mistral-7B and Qwen3-8B show moderate scale-related degradation.}
\label{tab:9gen}
\end{table}

The 9-generator LOGO yields 3-class F1 = 0.751 and binary F1 = 0.666 (Table~\ref{tab:9gen}).
Gemma-2-2B is a notable anomaly: gen-1 F1 = 0.036 indicates the classifier cannot distinguish Gemma's gen-1 output from gen-0.
This may reflect Gemma-2's architecture (grouped-query attention, RMSNorm) producing gen-1 distributional features unusually close to the seed text.
The method does not guarantee effectiveness across all generator architectures.
Larger models (Qwen3-8B at 0.647, Mistral-7B at 0.633) show moderate degradation, compressing feature margins as model capacity increases.

%==============================================================================
% 5. DISCUSSION
%==============================================================================
\section{Discussion}

\paragraph{Scope: base models only.}
Instruction-tuned models show substantially degraded depth signal (Table~\ref{tab:instruct_app}).
OLMo-2-1B-Instruct drops to LOGO binary F1 = 0.338 (vs.\ 0.763 base), with 6/7 features exhibiting Spearman $\rho$ direction reversal.
Llama-3.2-3B-Instruct is minimally affected ($-$1.1 pp), while Qwen3-1.7B-Instruct shows intermediate degradation ($-$16.3 pp).
This heterogeneity means we cannot predict \textit{a priori} which instruct models will degrade.
Base models remain relevant for data contamination analysis, as web-scraped synthetic text is predominantly generated by base or API models.

\paragraph{Temperature sensitivity.}
Depth signal is reliable within $t \in [0.5, 0.9)$ but degrades severely at $t = 0.9$ (binary F1 drops to 0.480; see Table~\ref{tab:temp_app}).
High-temperature sampling inflates entropy uniformly across depths, compressing inter-depth differences.
Cross-temperature transfer is asymmetric: training at $t = 1.0$ and testing at $t = 0.7$ yields F1 = 0.944, but the reverse yields only 0.418.

\paragraph{The \texttt{entropy\_std} question.}
Bootstrap TOST validation (1000 resamples, margin $d = 0.2$) rejects stationarity for three shape features (rejection rate 1000/1000) but \emph{never} rejects for \texttt{entropy\_std} (rate 0/1000, median $d = 0.458$; Table~\ref{tab:tost_app}).
Cross-reference analysis shows $d$ ranging from 0.310 (OLMo) to 0.520 (Qwen), further demonstrating inconsistency.
We adopt a spectrum model: drifting ($|\rho| > 0.50$) $\to$ intermediate (0.15--0.50) $\to$ stationary ($< 0.15$), with \texttt{entropy\_std} and \texttt{entropy\_min} in the intermediate category.

\paragraph{Handcrafted vs.\ end-to-end.}
QLoRA fine-tuned classifiers substantially outperform the 7-feature ensemble (Table~\ref{tab:qlora_app}): Qwen3-1.7B-Instruct achieves 3-class F1 = 0.907 and binary F1 = 0.867, a $+$14.5 pp improvement.
This confirms that depth signal is richer than what 7 features capture.
We position the 7-feature ensemble on \emph{interpretability and efficiency}: it requires only a 410M reference model and XGBoost ($\sim$150$\times$ less inference cost), and its features enable the DDR analysis that constitutes our primary conceptual contribution.

\paragraph{Semantic convergence vs.\ distributional divergence.}
Topic drift analysis reveals a counter-intuitive pattern: deeper generations converge in semantic space (cosine similarity increases from 0.622 at gen-0$\to$gen-1 to 0.859 at gen-3$\to$gen-4) while distributional features diverge.
DDR is a distributional phenomenon, not a semantic one.

%==============================================================================
% 6. CONCLUSION
%==============================================================================
\section{Conclusion}

We formulate generation depth classification for non-paraphrastic LLM generation chains and characterize differential depth response across 8 base-model generators from 5 families.
A lightweight 7-feature ensemble achieves leak-free binary F1 = 0.694 (per-generator range 0.633--0.724) under LOGO evaluation, and 0.830 on cross-model chains, using a single 410M reference model.
DDR is generation-intrinsic: it persists across reference models and is driven by drifting central-tendency features ($|\rho| > 0.60$) while shape features remain stationary ($|\rho| < 0.10$).
Depth-aware data curation does not improve downstream CLM quality over random subsampling, establishing depth fingerprinting as diagnostic rather than curative.

Current limitations include restriction to base models (instruct tuning degrades or reverses DDR), temperature sensitivity ($t \geq 0.9$), a 3-class resolution ceiling, and architecture-specific failures (Gemma-2-2B gen-1 F1 = 0.036).
Extending to instruct models and larger generators ($>$8B) remains future work.

%==============================================================================
% LIMITATIONS
%==============================================================================
\section*{Limitations}

\begin{enumerate}
    \item \textbf{Base models only.} The method does not reliably apply to instruction-tuned models, which represent a large fraction of deployed LLMs.
    \item \textbf{Temperature range.} Reliable only for $t \in [0.5, 0.9)$; most API defaults use $t = 1.0$.
    \item \textbf{Scale.} Larger generators ($>$3B) show moderate signal degradation; 70B+ models are untested.
    \item \textbf{Architecture sensitivity.} Gemma-2-2B represents a failure case; we cannot predict which architectures will fail.
    \item \textbf{Downstream evidence.} The negative finding is based on 3 seeds and a single perplexity metric, though seed variance is extremely small ($<$0.002).
    \item \textbf{Prompt-type dependence.} Performance varies across continuation, elaboration, and rewriting prompts (5 pp range).
    \item \textbf{Paraphrastic chains.} The method does not apply to paraphrastic chains (binary F1 = 0.55 vs.\ 0.694 for continuation).
\end{enumerate}

%==============================================================================
% REFERENCES
%==============================================================================
\bibliographystyle{plainnat}
\begin{thebibliography}{20}

\bibitem[Mitchell et al.(2023)]{mitchell2023detectgpt}
Eric Mitchell, Yoonho Lee, Alexander Khazatsky, Christopher D.\ Manning, and Chelsea Finn. 2023.
\newblock {DetectGPT}: Zero-Shot Machine-Generated Text Detection using Probability Curvature.
\newblock In \emph{Proceedings of ICML}.

\bibitem[Bao et al.(2024)]{bao2024fast}
Guangsheng Bao, Yanbin Zhao, Zhiyang Teng, Linyi Yang, and Yue Zhang. 2024.
\newblock Fast-{DetectGPT}: Efficient Zero-Shot Detection of Machine-Generated Text via Conditional Probability Curvature.
\newblock In \emph{Proceedings of ICLR}.

\bibitem[Gehrmann et al.(2019)]{gehrmann2019gltr}
Sebastian Gehrmann, Hendrik Strobelt, and Alexander Rush. 2019.
\newblock {GLTR}: Statistical Detection and Visualization of Generated Text.
\newblock In \emph{Proceedings of ACL: System Demonstrations}.

\bibitem[Hans et al.(2024)]{hans2024binoculars}
Abhimanyu Hans, Avi Schwarzschild, Valeriia Cherepanova, Hamid Kazemi, Aniruddha Saha, Micah Goldblum, Jonas Geiping, and Tom Goldstein. 2024.
\newblock Spotting {LLMs} With Binoculars: Zero-Shot Detection of Machine-Generated Text.
\newblock In \emph{Proceedings of ICML}.

\bibitem[Yang et al.(2024)]{yang2024dnagpt}
Xianjun Yang, Wei Cheng, Yue Wu, Linda Petzold, William Yang Wang, and Haifeng Chen. 2024.
\newblock {DNA-GPT}: Divergent N-Gram Analysis for Training-Free Detection of {GPT}-Generated Text.
\newblock In \emph{Proceedings of ICLR}.

\end{thebibliography}

%==============================================================================
% APPENDIX
%==============================================================================
\clearpage
\appendix

\section{Prompt Type Comparison}
\label{app:prompt}

Three prompt types---continuation, elaboration, and inspired rewriting---exhibit different DDR strengths (Table~\ref{tab:prompt}).
Continuation shows the strongest depth signal (CV F1 = 0.781), consistent with its largest semantic divergence per step.
Cross-prompt-type transfer (LOPTO) averages F1 = 0.753.

\begin{table}[h]
\centering
\small
\begin{tabular}{lcccc}
\toprule
\textbf{Prompt} & \textbf{CV F1} & \textbf{LOPTO F1} & \textbf{Ret.} & \textbf{$\rho_\text{hapax}$} \\
\midrule
Continuation & 0.781 & 0.762 & 0.831 & $-0.74$ \\
Rewriting & 0.737 & 0.733 & 0.939 & $-0.80$ \\
Elaboration & 0.730 & 0.762 & 0.986 & $-0.83$ \\
\bottomrule
\end{tabular}
\caption{Prompt type comparison. Higher semantic divergence (lower retention) correlates with stronger depth signal. LOPTO: leave-one-prompt-type-out.}
\label{tab:prompt}
\end{table}

\section{Paraphrastic vs.\ Non-Paraphrastic Chains}
\label{app:para}

Paraphrastic chains show severely weakened depth signal (Table~\ref{tab:para}), confirming the method's scope boundary.
Paraphrasing preserves semantics (retention = 0.97), compressing distributional variation across depths.
Binary F1 drops to 0.55, near chance.

\begin{table}[h]
\centering
\small
\begin{tabular}{lcccc}
\toprule
\textbf{Chain Type} & \textbf{CV F1} & \textbf{Bin. F1} & \textbf{$\rho_\text{hapax}$} & \textbf{Ret.} \\
\midrule
Continuation & 0.781 & 0.694 & $-0.74$ & 0.831 \\
Paraphrastic & 0.620 & 0.550 & $-0.35$ & 0.970 \\
$\Delta$ & $-0.161$ & $-0.144$ & $+0.39$ & $+0.139$ \\
\bottomrule
\end{tabular}
\caption{Paraphrastic vs.\ continuation chains. Semantic preservation compresses depth signal.}
\label{tab:para}
\end{table}

\section{Multi-Domain Analysis}
\label{app:domain}

DDR trends are preserved across four domains (Table~\ref{tab:domain}), though classification performance is domain-dependent.
Academic text shows the weakest signal (0.628), likely due to inherently low vocabulary diversity in specialized scientific writing.

\begin{table}[h]
\centering
\small
\begin{tabular}{lccc}
\toprule
\textbf{Domain} & \textbf{Binary F1} & \textbf{95\% CI} & \textbf{$\rho_\text{hapax}$} \\
\midrule
News & 0.712 & [0.690, 0.733] & $-0.65$ \\
Wikipedia & 0.695 & [0.675, 0.715] & $-0.63$ \\
Code & 0.687 & [0.664, 0.710] & $-0.58$ \\
Academic & 0.628 & [0.603, 0.652] & $-0.52$ \\
\bottomrule
\end{tabular}
\caption{Multi-domain binary F1. Cross-domain transfer F1 = 0.650.}
\label{tab:domain}
\end{table}

\section{Finer-Grained Ordinal Classification}
\label{app:ordinal}

5-class ordinal classification drops to F1 = 0.50 (Table~\ref{tab:ordinal}).
Ordinal methods (CORAL, Frank-Hall) provide no improvement, confirming the bottleneck is in feature signal, not model structure.
Gen-2 and gen-3 concentrate 65\% of confusion errors.

\begin{table}[h]
\centering
\small
\begin{tabular}{llccl}
\toprule
\textbf{$K$} & \textbf{Method} & \textbf{F1} & \textbf{G0 F1} & \textbf{G2+ F1} \\
\midrule
3 & XGBoost & 0.762 & 0.95 & 0.66 \\
5 & XGBoost & 0.500 & 0.92 & .35/.30/.38 \\
5 & CORAL & 0.500 & 0.91 & .34/.31/.38 \\
5 & Frank-Hall & 0.490 & 0.90 & .33/.30/.37 \\
\bottomrule
\end{tabular}
\caption{3-class vs.\ 5-class classification. 3-class is the practical resolution ceiling.}
\label{tab:ordinal}
\end{table}

\section{Robustness to Human Editing}
\label{app:edit}

The method shows graceful degradation under simulated human editing (Table~\ref{tab:edit}).
At 5\% random word replacement, binary F1 drops 4.6 pp; LLM rephrasing ($\sim$50\% edit distance) causes a comparable 4.5 pp drop.
Heavy editing ($>$20\%) substantially degrades signal.

\begin{table}[h]
\centering
\small
\begin{tabular}{lccc}
\toprule
\textbf{Edit Type} & \textbf{Rate} & \textbf{Bin. F1} & \textbf{$\Delta$} \\
\midrule
None & 0\% & 0.695 & — \\
Random & 5\% & 0.650 & $-0.046$ \\
Random & 10\% & 0.614 & $-0.081$ \\
Random & 20\% & 0.591 & $-0.104$ \\
LLM rephrase & $\sim$50\% ed. & 0.647 & $-0.045$ \\
\bottomrule
\end{tabular}
\caption{Robustness to simulated human editing.}
\label{tab:edit}
\end{table}

\section{LD-Score Feature Addition}
\label{app:ld}

Adding LD-Score (character-level distribution signature) as an 8th feature yields marginal improvement ($+$0.7 pp macro F1, McNemar $p = 0.12$; Table~\ref{tab:ld}).
LD-Score correlates highly with hapax\_ratio ($r = 0.68$), confirming character-level and token-level vocabulary diversity are largely redundant.

\begin{table}[h]
\centering
\small
\begin{tabular}{lccc}
\toprule
\textbf{Feature Set} & \textbf{$N$} & \textbf{Macro F1} & \textbf{McNemar $p$} \\
\midrule
7-feature & 7 & 0.804 & — \\
+ LD-Score & 8 & 0.811 & 0.12 \\
LD-Score only & 1 & 0.721 & — \\
\bottomrule
\end{tabular}
\caption{LD-Score feature addition. Marginal contribution is not significant.}
\label{tab:ld}
\end{table}

\section{End-to-End QLoRA Classification}
\label{app:qlora}

QLoRA fine-tuned classifiers substantially outperform the 7-feature ensemble (Table~\ref{tab:qlora_app}), confirming depth signal is richer than captured by 7 features.

\begin{table}[h]
\centering
\small
\begin{tabular}{lcccc}
\toprule
\textbf{Method} & \textbf{3-Cl F1} & \textbf{Bin F1} & \textbf{5-Cl F1} & \textbf{Cost} \\
\midrule
7-feat XGBoost & 0.762 & 0.694 & 0.500 & 1$\times$ \\
QLoRA OLMo-1B-Inst & 0.889 & 0.836 & 0.598 & $\sim$150$\times$ \\
QLoRA Qwen3-1.7B-Inst & 0.907 & 0.867 & 0.551 & $\sim$250$\times$ \\
QLoRA OLMo-1B-Base & 0.859 & 0.793 & 0.540 & $\sim$150$\times$ \\
QLoRA Qwen3-1.7B-Base & 0.904 & 0.860 & 0.554 & $\sim$250$\times$ \\
\bottomrule
\end{tabular}
\caption{Handcrafted vs.\ end-to-end classification. QLoRA achieves $+$12--15 pp improvement at $\sim$150--250$\times$ inference cost.}
\label{tab:qlora_app}
\end{table}

\section{Text Length Control}
\label{app:length}

Length-depth correlation is weak ($|\rho| < 0.14$; Table~\ref{tab:length}).
Length-matched subsampling reduces binary F1 by 4.6 pp, but a random-subsample control (same $N$) also drops 3.0 pp, attributing only 1.6 pp to the length confound.

\begin{table}[h]
\centering
\small
\begin{tabular}{lccc}
\toprule
\textbf{Condition} & \textbf{Bin. F1} & \textbf{$\Delta$} & \textbf{Note} \\
\midrule
Full (unmatched) & 0.770 & — & Full dataset \\
Length-matched & 0.724 & $-0.046$ & Equalized \\
Random subsample & 0.740 & $-0.030$ & Size control \\
$\Delta$ (length) & — & $-0.016$ & Attributable \\
\bottomrule
\end{tabular}
\caption{Length confound analysis. Only 1.6 pp of the 4.6 pp drop is attributable to length.}
\label{tab:length}
\end{table}

\section{Large-Scale Generator}
\label{app:scale}

Adding Qwen3-8B-Base (8B parameters) confirms moderate scale degradation (Table~\ref{tab:scale}).

\begin{table}[h]
\centering
\small
\begin{tabular}{lcccc}
\toprule
\textbf{Generator} & \textbf{Params} & \textbf{LOGO Bin F1} & \textbf{$\rho_\text{hapax}$} & \textbf{$\rho_\text{ent.m}$} \\
\midrule
Pythia-1B & 1B & 0.718 & $-0.65$ & $-0.64$ \\
Qwen3-0.6B & 0.6B & 0.659 & $-0.60$ & $-0.58$ \\
Pythia-2.8B & 2.8B & 0.724 & $-0.63$ & $-0.60$ \\
Qwen3-1.7B & 1.7B & 0.702 & $-0.61$ & $-0.57$ \\
OLMo-2-1B & 1.5B & 0.668 & $-0.57$ & $-0.52$ \\
Qwen3-8B & 8B & 0.683 & $-0.56$ & $-0.48$ \\
\midrule
\textbf{6-gen avg} & — & \textbf{0.692} & $-0.60$ & $-0.57$ \\
\bottomrule
\end{tabular}
\caption{6-generator LOGO including Qwen3-8B. Scale degradation is moderate ($-$1.1 pp average).}
\label{tab:scale}
\end{table}

\section{DDR Causality Analysis}
\label{app:causality}

Three independent methods confirm drifting features carry $\sim$5$\times$ more predictive signal than stationary features (Table~\ref{tab:causality}).

\begin{table}[h]
\centering
\small
\begin{tabular}{lccc}
\toprule
\textbf{Method} & \textbf{Drifting} & \textbf{Stationary} & \textbf{Ratio} \\
\midrule
Perm.\ importance & 0.158 & 0.030 & 5.3$\times$ \\
Granger causality & 6/7 sig & — & — \\
Feature intervention & 0.020 & 0.006 & 3.3$\times$ \\
\bottomrule
\end{tabular}
\caption{DDR causality evidence. Drifting features dominate across all methods.}
\label{tab:causality}
\end{table}

\section{TOST Bootstrap Stability}
\label{sec:tost}

Bootstrap TOST (1000 resamples, $d = 0.2$ margin) validates stationarity for entropy\_max, entropy\_min, and tail\_decay (Table~\ref{tab:tost_app}) but rejects it for \texttt{entropy\_std} (rate 0/1000, median $d = 0.458$).

\begin{table}[h]
\centering
\small
\begin{tabular}{lcccc}
\toprule
\textbf{Feature} & \textbf{Med.\ $d$} & \textbf{95\% CI} & \textbf{Rej.\ Rate} & \textbf{Class} \\
\midrule
entropy\_max & 0.062 & [0.037, 0.088] & 1000/1000 & Stationary \\
entropy\_min & $-0.016$ & [$-0.042$, 0.011] & 1000/1000 & Stationary \\
tail\_decay & 0.088 & [0.065, 0.111] & 1000/1000 & Stationary \\
\textbf{entropy\_std} & \textbf{0.458} & \textbf{[0.435, 0.479]} & \textbf{0/1000} & \textbf{Intermediate} \\
\bottomrule
\end{tabular}
\caption{TOST bootstrap stability. \texttt{entropy\_std} fails stationarity at $d = 0.2$ margin.}
\label{tab:tost_app}
\end{table}

\section{Instruct Model Evaluation}
\label{app:instruct}

Instruction tuning substantially degrades depth signal, with heterogeneous effects across models (Table~\ref{tab:instruct_app}).

\begin{table}[h]
\centering
\small
\begin{tabular}{lcccc}
\toprule
\textbf{Model} & \textbf{Base F1} & \textbf{Inst F1} & \textbf{$\Delta$} & \textbf{Feat Rev.} \\
\midrule
OLMo-2-1B & 0.763 & 0.338 & $-0.425$ & 6/7 \\
Qwen3-1.7B & 0.821 & 0.658 & $-0.163$ & partial \\
Llama-3.2-3B & 0.849 & 0.838 & $-0.011$ & 3/7 \\
\bottomrule
\end{tabular}
\caption{Instruct vs.\ base model depth signal. OLMo-Instruct shows near-complete degradation with feature direction reversal.}
\label{tab:instruct_app}
\end{table}

\section{Temperature Sensitivity}
\label{app:temp}

\begin{table}[h]
\centering
\small
\begin{tabular}{lccc}
\toprule
\textbf{$t$} & \textbf{Within-CV Bin F1} & \textbf{$\rho_\text{ent.m}$} & \textbf{Reliable} \\
\midrule
0.5 & 0.813 & $-0.58$ & Yes \\
0.7 & 0.694 & $-0.60$ & Yes \\
0.9 & 0.480 & $-0.30$ & No \\
\bottomrule
\end{tabular}
\caption{Temperature sensitivity. Signal degrades severely at $t \geq 0.9$.}
\label{tab:temp_app}
\end{table}

Cross-temperature transfer (Qwen3-1.7B): training at $t = 0.7$, testing at $t = 1.0$ yields binary F1 = 0.418 ($-$37.1 pp vs.\ within-temperature). The reverse direction ($t = 1.0 \to 0.7$) yields 0.944 ($+$15.5 pp), demonstrating asymmetric transfer.

\section{Independent Seed Validation}
\label{app:seed_val}

Independent seed sets confirm shared-seed LOGO does not suffer seed leakage (Table~\ref{tab:seed_val}).

\begin{table}[h]
\centering
\small
\begin{tabular}{lcc}
\toprule
\textbf{Setup} & \textbf{Binary F1} & \textbf{Seeds Shared} \\
\midrule
Shared-seed LOGO & 0.694 & Yes \\
Independent: Pythia$\to$OLMo & 0.695 & No \\
Independent: OLMo$\to$Pythia & 0.666 & No \\
Independent avg & 0.680 & No \\
\bottomrule
\end{tabular}
\caption{Seed leakage validation. $\Delta = -0.014$, within noise.}
\label{tab:seed_val}
\end{table}

\section{Seed Complexity Stratification}
\label{app:seed_complex}

Depth signal is independent of seed text complexity (Table~\ref{tab:seed_complex}; max difference 1.06 pp).

\begin{table}[h]
\centering
\small
\begin{tabular}{lccc}
\toprule
\textbf{Complexity} & \textbf{Bin F1} & \textbf{Multi F1} & \textbf{$N$} \\
\midrule
Low & 0.685 & 0.772 & $\sim$332 \\
Medium & 0.695 & 0.778 & $\sim$332 \\
High & 0.695 & 0.779 & $\sim$331 \\
\bottomrule
\end{tabular}
\caption{Seed complexity stratification. Depth signal is seed-independent.}
\label{tab:seed_complex}
\end{table}

\section{Topic Drift Correlation}
\label{app:topic}

Semantic similarity increases with depth (Table~\ref{tab:topic}), while distributional features diverge---confirming DDR is distributional, not semantic.

\begin{table}[h]
\centering
\small
\begin{tabular}{lcc}
\toprule
\textbf{Transition} & \textbf{Mean Cosine Sim} & \textbf{Std} \\
\midrule
gen-0 $\to$ gen-1 & 0.622 & 0.165 \\
gen-1 $\to$ gen-2 & 0.683 & 0.184 \\
gen-2 $\to$ gen-3 & 0.801 & 0.201 \\
gen-3 $\to$ gen-4 & 0.859 & 0.190 \\
\bottomrule
\end{tabular}
\caption{Semantic similarity by depth transition. Deeper generations converge semantically while diverging distributionally.}
\label{tab:topic}
\end{table}

\section{TOST Equivalence Tests (Downstream)}
\label{app:tost_downstream}

\begin{table}[h]
\centering
\small
\begin{tabular}{lccc}
\toprule
\textbf{Comparison} & \textbf{Margin} & \textbf{$p$} & \textbf{Equiv.} \\
\midrule
SM vs.\ depth01 & $\pm$2\% & $1.1 \times 10^{-4}$ & Yes \\
SM vs.\ binary-filt. & $\pm$2\% & $8.3 \times 10^{-5}$ & Yes \\
All vs.\ SM & $\pm$2\% & 0.15 & No \\
All vs.\ SM & $\pm$3\% & 0.003 & Yes \\
\bottomrule
\end{tabular}
\caption{TOST equivalence tests for downstream CLM quality. SM = size-matched random.}
\label{tab:tost_downstream}
\end{table}

\section{\texttt{entropy\_std} Cross-Reference Inconsistency}
\label{app:entropy_std}

\begin{table}[h]
\centering
\small
\begin{tabular}{lccc}
\toprule
\textbf{Ref.\ Model} & \textbf{$d$} & \textbf{TOST Rej.} & \textbf{Class} \\
\midrule
Pythia-410M & 0.458 & 0/1000 & Contested \\
OLMo-2-1B & 0.310 & partial & Borderline \\
Qwen3-0.6B & 0.520 & 0/1000 & Contested \\
\bottomrule
\end{tabular}
\caption{\texttt{entropy\_std} Cohen's $d$ across reference models.}
\label{tab:entropy_std_cross}
\end{table}

\end{document}
